
\section{\BookRecordIndex{T0220}拉普拉斯算子的\BookRecordIndex{T0408}正则性}\label{sec:B-3}

现在可以建立\BookRecordIndex{T0220}拉普拉斯算子的关键\BookRecordIndex{T0408}正则性定理与估计。

\begin{theorem}[内部\BookRecordIndex{T0408}正则性]\label{thm:B-3-1}
设 $1<p<\infty$，$k\ge 0$ 为整数，$\Omega\subset\mathbb R^n$ 为开域。
若 $u\in L^1_{\mathrm{loc}}(\Omega)$ 是方程
\[
  \Delta u=f,\qquad f\in W^{k,p}_{\mathrm{loc}}(\Omega)
\]
的\BookRecordIndex{T0295}弱解，则 $u\in W^{k+2,p}_{\mathrm{loc}}(\Omega)$。此外，对每个有界开集
$\Omega'\subset\mathbb R^n$，若 $\overline{\Omega'}\subset\Omega$，则存在常数
$c=c(k,p,n,\Omega',\Omega)>0$，使每个
$u\in C^\infty(\overline\Omega)$ 都满足
\[
  \|u\|_{W^{k+2,p}(\Omega')}
  \le c\bigl(\|\Delta u\|_{W^{k,p}(\Omega)}
  +\|u\|_{L^p(\Omega)}\bigr).
\]
\end{theorem}

\begin{proof}
设 $u,f,\Omega'$ 如定理所述。取包含 $\overline{\Omega'}$ 的有界开邻域
$U$，使 $\overline U\subset\Omega$。当 $n=1$ 时，
$u''=f$ 可在各区间上直接积分，以下\BookRecordIndex{T0205}卷积论证取 $n\ge2$。
\sourcefootnote{原书在一般 $n$ 下直接使用 $K$，但其定义式~\eqref{eq:B-2-1}
只覆盖 $n\ge2$。一维情形可由 $u''=f$ 直接处理；下文定理 B.3.2 的
\BookRecordIndex{T0205}卷积证明有同一省略。}
再取光滑\BookRecordIndex{T0188}截断函数
$\beta\in C^\infty_0(\Omega)$，使 $\beta|_U=1$，并定义
\[
  v:=K*(\beta f).
\]
我们证明 $v\in W^{k+2,p}_{\mathrm{loc}}(\Omega)$。为此，取一列支集同处
$\Omega$ 中某个固定紧集内的 $f_i\in C^\infty_0(\Omega)$，使
\[
  \lim_{i\to\infty}\|f_i-\beta f\|_{W^{k,p}(\Omega)}=0.
\]
令 $v_i:=K*f_i$。每个 $v_i$ 都是光滑的，因为局部可积函数与紧支光滑函数
的\BookRecordIndex{T0205}卷积是光滑函数。对任意有界开集 $V$，若
$\overline U\subset V$ 且 $\overline V\subset\Omega$，
$K$ 的局部可积性与 Young 不等式给出
$v_i\to v$ 于 $L^p(V)$；定理~\ref{thm:B-2-7} 还表明
$v_i$ 于 $W^{k+2,p}(V)$ 中为柯西列。因此
$v\in W^{k+2,p}(V)$。
\sourcefootnote{原书此处把 $L^p$ 收敛、$W^{k+2,p}$ 柯西性以及结论
均写在整个 $\Omega$ 上。但 $\Omega$ 未假设有界；例如在二维，
$K*\beta f$ 可在无穷远处呈对数增长，不必属于 $L^p(\Omega)$。
证明所需的结论是在 $U$ 附近的局部\BookRecordIndex{T0408}正则性，故此处将这些断言限制在
相对紧的有界开集 $V$ 上。}
由引理~\ref{lem:B-2-2} 的正向证明及逼近
\sourcefootnote{原书此处直接引用引理 B.2.2；该引理按其陈述要求
解与右端均有紧支撑，而这里的 $v=K*(\beta f)$ 一般没有紧支撑。
所需的正向\BookRecordIndex{T0123}分布恒等式可由该引理证明中的计算经逼近得到。}，
$v$ 是 $\Delta v=\beta f$ 的\BookRecordIndex{T0295}弱解。
于是 $u-v$ 在 $U$ 上弱满足 $\Delta(u-v)=0$。
由 Weyl 引理，$u-v$ 在 $U$ 上实解析，故
$u\in W^{k+2,p}(\Omega')$。第一个断言得证。

为证明估计，先注意：对每个有界开集 $\Omega\subset\mathbb R^n$ 及每个整数
$k\ge0$，存在常数 $c=c(k,n,p,\Omega)$，使所有
$u\in C^\infty_0(\Omega)$ 满足
\begin{equation}\label{eq:B-3-1}
  \|u\|_{W^{k+2,p}(\Omega)}
  \le c\|\Delta u\|_{W^{k,p}(\Omega)}.
\end{equation}
当 $k=0$ 时，这由推论~\ref{cor:B-2-8} 和引理~\ref{lem:B-1-6} 得出；
$k\ge1$ 时，可用归纳法证明（见习题 B.1.2）。由于
$C^\infty_0(\Omega)$ 在 $W^{k+2,p}_0(\Omega)$ 中稠密，
不等式~\eqref{eq:B-3-1} 也对每个
$u\in W^{k+2,p}_0(\Omega)$ 成立。

其次证明：对每个满足 $\overline{\Omega'}\subset\Omega$ 的有界开集
$\Omega'$
\sourcefootnote{原书在此估计及下文定理 B.3.2 中只写
$\overline{\Omega'}\subset\Omega$，未要求 $\Omega'$ 有界。
然而其证明要在 $\Omega'$ 外选支于 $\Omega$ 内的\BookRecordIndex{T0188}截断函数；
对一般无界 $\Omega'$，这种\BookRecordIndex{T0188}截断函数未必存在。这里补入
$\Omega'$ 有界这一充分条件。}
和每个 $p>1$，存在常数 $c=c(n,p,\Omega',\Omega)>0$，
使每个 $u\in W^{2,p}_{\mathrm{loc}}(\Omega)$ 满足
\begin{equation}\label{eq:B-3-2}
  \|u\|_{W^{1,p}(\Omega')}
  \le c\bigl(\|\Delta u\|_{L^p(\Omega)}
  +\|u\|_{L^p(\Omega)}\bigr).
\end{equation}
取 $\beta\in C^\infty_0(\Omega)$，使它在
$\overline{\Omega'}$ 的某个有界开邻域 $U$ 上等于 $1$，并令
$v:=K*(\beta\Delta u)$。由于 $\nabla v=\nabla K*(\beta\Delta u)$，
且 $\nabla K$ 在紧集上可积，Young 不等式给出以下局部估计
\sourcefootnote{原书这条\BookRecordIndex{T0205}卷积估计的左端写为
$\|v\|_{W^{1,p}(\Omega)}$。当 $\Omega$ 无界时，全局范数仍可能不存在；
在证明中只须取有界的 $U$。}：
\[
  \|v\|_{W^{1,p}(U)}
  \le c_1\|\beta\Delta u\|_{L^p(\Omega)}
  \le c_1\|\Delta u\|_{L^p(\Omega)}.
\]
函数 $v-u$ 在 $U$ 上调和。于是，\BookRecordIndex{T0104}调和函数及其一阶导数的\BookRecordIndex{T0267}平均值性质表明，
存在常数 $c_2>0$，使
\[
  \|v-u\|_{W^{1,p}(\Omega')}
  \le c_2\|v-u\|_{L^p(U)}
  \le c_2\bigl(\|v\|_{L^p(U)}+\|u\|_{L^p(U)}\bigr).
\]
将这些不等式相加，即得~\eqref{eq:B-3-2}。

现在设 $u\in W^{k+2,p}_{\mathrm{loc}}(\Omega)$，记
$f:=\Delta u\in W^{k,p}_{\mathrm{loc}}(\Omega)$，并设
$\overline{\Omega'}\subset\Omega$。取有界开集
$\Omega_0,\ldots,\Omega_{k+1}$，使
\[
  \overline{\Omega'}\subset\Omega_{k+1},\qquad
  \overline{\Omega_{j+1}}\subset\Omega_j,\qquad
  \overline{\Omega_0}\subset\Omega.
\]
对 $j=0,\ldots,k$，取
$\beta_j\in C^\infty_0(\Omega_j)$，使
$\beta_j|_{\Omega_{j+1}}=1$。于是
$\beta_j u\in W^{k+2,p}_0(\Omega_j)$，由~\eqref{eq:B-3-1} 得
\begin{align*}
  \|u\|_{W^{j+2,p}(\Omega_{j+1})}
  &\le \|\beta_j u\|_{W^{j+2,p}(\Omega_j)}\\
  &\le c_j\|\Delta(\beta_j u)\|_{W^{j,p}(\Omega_j)}\\
  &\le c'_j\bigl(\|\Delta u\|_{W^{j,p}(\Omega_j)}
    +\|u\|_{W^{j+1,p}(\Omega_j)}\bigr),
  \qquad j=0,\ldots,k.
\end{align*}
归纳可得
\[
  \|u\|_{W^{k+2,p}(\Omega_{k+1})}
  \le c\bigl(\|\Delta u\|_{W^{k,p}(\Omega_0)}
  +\|u\|_{W^{1,p}(\Omega_0)}\bigr).
\]
再将~\eqref{eq:B-3-2} 中的 $\Omega'$ 换成 $\Omega_0$，即得所需估计。
\end{proof}

有时需要考虑 $\Delta u=f$ 的\BookRecordIndex{T0295}弱解，其中 $f$ 不是函数，而是
$W^{-1,p}$ 中的\BookRecordIndex{T0123}分布。下面用方程
\[
  \Delta u=\operatorname{div}f,
  \qquad f\in L^p
\]
的\BookRecordIndex{T0295}弱解来表述这种情形。

\begin{theorem}\label{thm:B-3-2}
设 $1<p<\infty$，$\Omega'\subset\Omega\subset\mathbb R^n$ 为开集，
其中 $\Omega'$ 有界
\sourcefootnote{原书此定理未要求 $\Omega'$ 有界；
此处补入该条件，以保证证明中可取紧支于 $\Omega$、
在 $\Omega'$ 的邻域上为 $1$ 的\BookRecordIndex{T0188}截断函数。}，
且 $\overline{\Omega'}\subset\Omega$。则存在常数
$c=c(p,n,\Omega',\Omega)>0$，使下述结论成立。
设 $u\in L^1_{\mathrm{loc}}(\Omega)$，
$f=(f_0,\ldots,f_n)\in L^p_{\mathrm{loc}}(\Omega,\mathbb R^{n+1})$，
且对每个 $\phi\in C^\infty_0(\Omega)$ 有
\[
  \int_\Omega u(x)\Delta\phi(x)\,dx
  =\int_\Omega f_0(x)\phi(x)\,dx
   -\sum_{j=1}^n\int_\Omega f_j(x)\partial_j\phi(x)\,dx.
\]
则 $u\in W^{1,p}_{\mathrm{loc}}(\Omega)$，并且
\[
  \|u\|_{W^{1,p}(\Omega')}
  \le c\bigl(\|f\|_{L^p(\Omega)}+\|u\|_{L^p(\Omega)}\bigr).
\]
\end{theorem}

\begin{proof}
当 $n=1$ 时，由 $u''=f_0+\partial_1f_1$ 在区间上直接积分即可；
以下设 $n\ge2$。
取 $\overline{\Omega'}$ 的有界开邻域 $U$，使
$\overline U\subset\Omega$。令 $\beta\in C^\infty_0(\Omega)$
为在 $U$ 上等于 $1$ 的光滑\BookRecordIndex{T0188}截断函数，并定义
\[
  v:=K*(\beta f_0)+\sum_{j=1}^n K_j*(\beta f_j).
\]
定理~\ref{thm:B-2-7} 以及与定理~\ref{thm:B-3-1} 的证明相同的逼近论证
（留给读者）表明 $v\in W^{1,p}_{\mathrm{loc}}(\Omega)$，且有估计
\[
  \|v\|_{W^{1,p}(U)}
  \le c_1\|\beta f\|_{L^p(\Omega)}
  \le c_1\|f\|_{L^p(\Omega)}.
\]
这里也用到了\BookRecordIndex{T0261}庞加莱不等式。由引理~\ref{lem:B-2-2} 的正向证明及逼近，
$v$ 是下式的\BookRecordIndex{T0295}弱解：
\[
  \Delta v=\beta f_0+\sum_{j=1}^n\partial_j(\beta f_j),
\]
其中 $\partial_j(\beta f_j)$ 理解为\BookRecordIndex{T0123}分布。因此 $u-v$ 在 $U$ 上弱满足
$\Delta(u-v)=0$。由 Weyl 引理，$u-v$ 在 $U$ 上调和，
从而 $u\in W^{1,p}(\Omega')$。又由\BookRecordIndex{T0104}调和函数的\BookRecordIndex{T0267}平均值性质，
存在常数 $c_2>0$，使
\[
  \|v-u\|_{W^{1,p}(\Omega')}
  \le c_2\|v-u\|_{L^p(U)}
  \le c_2\bigl(\|v\|_{W^{1,p}(U)}+\|u\|_{L^p(U)}\bigr).
\]
将这些不等式相加，即得所需估计。
\end{proof}

\begin{exercise}\label{ex:B-3-3}
设 $\Omega\subset\mathbb C$ 为开集。\BookRecordIndex{T0209}柯西--黎曼方程可写为
\[
  \frac{\partial w}{\partial\bar z}=h,
  \qquad
  \frac{\partial}{\partial\bar z}
  :=\frac12\left(\frac{\partial}{\partial s}
  +i\frac{\partial}{\partial t}\right),
\]
其中 $s+it$ 是 $\Omega$ 的坐标。为\BookRecordIndex{T0210}柯西--黎曼算子表述并证明
定理~\ref{thm:B-3-1} 的类似结果。

\emph{提示：}写 $w=u+iv$，$h=(f+ig)/2$。
证明：若 $w$ 是\BookRecordIndex{T0209}柯西--黎曼方程的（弱）解，则 $u,v$ 分别是二阶方程
\[
  \Delta u=\partial_s f+\partial_t g,
  \qquad
  \Delta v=\partial_s g-\partial_t f
\]
的（弱）解。或者证明
$N(z):=1/(\pi z)$ 是\BookRecordIndex{T0209}柯西--黎曼方程的\BookRecordIndex{T0173}基本解，其含义是：
任意两个函数 $w,h\in C^\infty_0(\Omega,\mathbb C)$ 满足
$\partial w/\partial\bar z=h$，当且仅当 $w=N*h$。
注意 $N=4\partial K/\partial z$，其中
$K(z)=(2\pi)^{-1}\log|z|$ 是拉普拉斯方程的\BookRecordIndex{T0173}基本解。
\end{exercise}
